% Source PDF page 11
\begin{frame}[t,allowframebreaks]{Commands}
  \fontsize{15}{18}\selectfont
  \begin{itemize}
    \item -A : Appends the \textcolor{CommandGreen}{\bfseries iptables} rule to the end of the specified
  \end{itemize}
  chain.\par
  \begin{itemize}
    \item -F : Flushes the selected chain, which effectively deletes every
  \end{itemize}
  rule in the chain. If no chain is specified, this command flushes every rule from every chain\par
  \begin{itemize}
    \item -L : Lists all of the rules in the chain specified after the
  \end{itemize}
  command. To list all rules in all chains in the default \textcolor{CommandGreen}{\bfseries filter} table, do not specify a chain or table. Otherwise, the following syntax should be used to list the rules in a specific chain in a particular table:\par
  \vspace{0.08cm}
  {\color{CommandGreen}\ttfamily\fontsize{13.5}{16.5}\selectfont iptables -L <chain-name> -t <table-name>\par}
  \begin{itemize}
    \item -N : Creates a new chain with a user-specified name.
    \item -P : Sets the default policy for a particular chain, so that when
  \end{itemize}
  packets traverse an entire chain without matching a rule, they will be sent on to a particular target, such as \textcolor{WarningColor}{\bfseries ACCEPT} or \textcolor{WarningColor}{\bfseries DROP}.\par
  \vspace{0.08cm}
  {\color{CommandGreen}\ttfamily\fontsize{13.5}{16.5}\selectfont iptables -P INPUT DROP\par}
\end{frame}
